% ID: SEL-D0072-OB-E01
% ESTADO: validada. S072.
\textbf{a)} Sean $x,y$ las cantidades de camisas A y B. El modelo consiste en maximizar $8x+6y$ euros con
\[
x+y\le100\,000,\qquad y\ge0{,}60(x+y),\qquad x,y\in\mathbb Z_{\ge0}.
\]
La segunda restricción equivale a $2y\ge3x$.

\textbf{b)} En el problema del recinto, la cota inferior es $(x-3)/4$ y las superiores son $x$ y $6-2x$. Al cruzar las fronteras por pares se obtienen $(-1,-1)$, $(2,2)$ y $(3,0)$, que cumplen las tres restricciones. El recinto es el triángulo cerrado.
\begin{center}
\setlength{\unitlength}{0.9mm}
\begin{picture}(103,57)(-31,-18)
\qbezier(-26,0)(18,0)(62,0)
\qbezier(0,-13)(0,11)(0,35)
\qbezier(-18,-11)(9,5.5)(36,22)
\qbezier(36,22)(45,11)(54,0)
\qbezier(54,0)(18,-5.5)(-18,-11)
\put(63,0){$x$}\put(2,35){$y$}
\put(-30,-15){\scriptsize $(-1,-1)$}
\put(34,25){\scriptsize $(2,2)$}
\put(54,3){\scriptsize $(3,0)$}
\end{picture}
\end{center}
Por la segunda restricción, $F=y+2x\le6$. La igualdad se alcanza en todo el lado entre $(2,2)$ y $(3,0)$. Por tanto, el máximo es $\boxed{6}$ en todo ese segmento.
